% Options for packages loaded elsewhere
\PassOptionsToPackage{unicode}{hyperref}
\PassOptionsToPackage{hyphens}{url}
\documentclass[
  10pt,
]{article}
\usepackage{xcolor}
\usepackage[margin=1in]{geometry}
\usepackage{amsmath,amssymb}
\setcounter{secnumdepth}{-\maxdimen} % remove section numbering
\usepackage{iftex}
\ifPDFTeX
  \usepackage[T1]{fontenc}
  \usepackage[utf8]{inputenc}
  \usepackage{textcomp} % provide euro and other symbols
\else % if luatex or xetex
  \usepackage{unicode-math} % this also loads fontspec
  \defaultfontfeatures{Scale=MatchLowercase}
  \defaultfontfeatures[\rmfamily]{Ligatures=TeX,Scale=1}
\fi
\usepackage{lmodern}
\ifPDFTeX\else
  % xetex/luatex font selection
\fi
% Use upquote if available, for straight quotes in verbatim environments
\IfFileExists{upquote.sty}{\usepackage{upquote}}{}
\IfFileExists{microtype.sty}{% use microtype if available
  \usepackage[]{microtype}
  \UseMicrotypeSet[protrusion]{basicmath} % disable protrusion for tt fonts
}{}
\makeatletter
\@ifundefined{KOMAClassName}{% if non-KOMA class
  \IfFileExists{parskip.sty}{%
    \usepackage{parskip}
  }{% else
    \setlength{\parindent}{0pt}
    \setlength{\parskip}{6pt plus 2pt minus 1pt}}
}{% if KOMA class
  \KOMAoptions{parskip=half}}
\makeatother
\usepackage{longtable,booktabs,array}
\usepackage{calc} % for calculating minipage widths
% Correct order of tables after \paragraph or \subparagraph
\usepackage{etoolbox}
\makeatletter
\patchcmd\longtable{\par}{\if@noskipsec\mbox{}\fi\par}{}{}
\makeatother
% Allow footnotes in longtable head/foot
\IfFileExists{footnotehyper.sty}{\usepackage{footnotehyper}}{\usepackage{footnote}}
\makesavenoteenv{longtable}
\usepackage{graphicx}
\makeatletter
\newsavebox\pandoc@box
\newcommand*\pandocbounded[1]{% scales image to fit in text height/width
  \sbox\pandoc@box{#1}%
  \Gscale@div\@tempa{\textheight}{\dimexpr\ht\pandoc@box+\dp\pandoc@box\relax}%
  \Gscale@div\@tempb{\linewidth}{\wd\pandoc@box}%
  \ifdim\@tempb\p@<\@tempa\p@\let\@tempa\@tempb\fi% select the smaller of both
  \ifdim\@tempa\p@<\p@\scalebox{\@tempa}{\usebox\pandoc@box}%
  \else\usebox{\pandoc@box}%
  \fi%
}
% Set default figure placement to htbp
\def\fps@figure{htbp}
\makeatother
% definitions for citeproc citations
\NewDocumentCommand\citeproctext{}{}
\NewDocumentCommand\citeproc{mm}{%
  \begingroup\def\citeproctext{#2}\cite{#1}\endgroup}
\makeatletter
 % allow citations to break across lines
 \let\@cite@ofmt\@firstofone
 % avoid brackets around text for \cite:
 \def\@biblabel#1{}
 \def\@cite#1#2{{#1\if@tempswa , #2\fi}}
\makeatother
\newlength{\cslhangindent}
\setlength{\cslhangindent}{1.5em}
\newlength{\csllabelwidth}
\setlength{\csllabelwidth}{3em}
\newenvironment{CSLReferences}[2] % #1 hanging-indent, #2 entry-spacing
 {\begin{list}{}{%
  \setlength{\itemindent}{0pt}
  \setlength{\leftmargin}{0pt}
  \setlength{\parsep}{0pt}
  % turn on hanging indent if param 1 is 1
  \ifodd #1
   \setlength{\leftmargin}{\cslhangindent}
   \setlength{\itemindent}{-1\cslhangindent}
  \fi
  % set entry spacing
  \setlength{\itemsep}{#2\baselineskip}}}
 {\end{list}}
\usepackage{calc}
\newcommand{\CSLBlock}[1]{\hfill\break\parbox[t]{\linewidth}{\strut\ignorespaces#1\strut}}
\newcommand{\CSLLeftMargin}[1]{\parbox[t]{\csllabelwidth}{\strut#1\strut}}
\newcommand{\CSLRightInline}[1]{\parbox[t]{\linewidth - \csllabelwidth}{\strut#1\strut}}
\newcommand{\CSLIndent}[1]{\hspace{\cslhangindent}#1}
\setlength{\emergencystretch}{3em} % prevent overfull lines
\providecommand{\tightlist}{%
  \setlength{\itemsep}{0pt}\setlength{\parskip}{0pt}}
\usepackage{bookmark}
\IfFileExists{xurl.sty}{\usepackage{xurl}}{} % add URL line breaks if available
\urlstyle{same}
\hypersetup{
  pdftitle={GapHarness: Obligation-First Minimal Harness Synthesis for API-Only LLM Agents},
  pdfauthor={Anonymous Authors},
  hidelinks,
  pdfcreator={LaTeX via pandoc}}

\title{GapHarness: Obligation-First Minimal Harness Synthesis for
API-Only LLM Agents}
\author{Anonymous Authors}
\date{}

\begin{document}
\maketitle

\section{Abstract}\label{abstract}

API-only LLM agents increasingly rely on external runtimes for
retrieval, code execution, state management, action, control, and
verification. Existing tool routers often choose tools directly, which
can conflate two separable questions: what external obligations a user
query imposes, and what smallest runtime harness can satisfy those
obligations. We introduce GapHarness, a research system that formulates
agent harnessing as obligation inference followed by minimal runtime
compilation. Given an obligation profile and a declared module registry,
GapHarness compiles the lowest-cost module subset that covers required
obligations and capabilities, executes a deterministic sandbox loop, and
verifies sufficiency and relative minimality.

We define an obligation ontology covering Observation, Execution, State,
Action, Control, and Verification, and introduce metrics for
over-harnessing, under-harnessing, wrong-harnessing, minimality regret,
and redundancy. On the 1000-task human-audited GapBench controlled
benchmark, GapHarness matches the oracle minimal harness under gold
obligations. With LLM-inferred obligations on held-out test800,
GapHarness reaches 0.89 success, outperforming direct and router
baselines while exposing a calibration gap. A registry-guarded profiler
reduces unsupported false positives from 56 to 12 and improves held-out
success to 0.94, at the cost of higher regret. Stress tests show that
success depends on declared registry affordances and meaningful
obligation labels; negative controls show that GapHarness avoids
tool-bait over-harnessing. Transfer runs on GAIA-Transfer and
GapBench-Natural identify boundary conditions rather than claiming full
open-world task solving.

\section{1. Introduction}\label{introduction}

LLM agents are commonly evaluated as end-to-end systems: a model
receives a request, selects tools, acts in an environment, and produces
an answer. This framing is useful, but it hides a smaller technical
problem that appears whenever an API-only model must decide what runtime
support is needed. A user request may require observation beyond the
prompt, deterministic execution, durable state, sandbox action, control
over risky operations, or independent verification. These are
obligations, not tools.

Tool-using and reasoning-action agents have shown that external actions
can improve model behavior (\citeproc{ref-yao2023react}{S. Yao et al.
2023}; \citeproc{ref-schick2023toolformer}{Schick et al. 2023};
\citeproc{ref-patil2023gorilla}{Patil et al. 2023};
\citeproc{ref-qin2024toolllm}{Qin et al. 2024}). However, choosing tools
is not the same as deciding which external obligations must be
satisfied. A keyword-based router may over-call tools when the user
explicitly requests no tools, and it may under-call tools when a
warranted answer requires evidence, state, or verification that is not
obvious from surface wording. Benchmarks of tool choice and tool
awareness also show that deciding whether and which tools to use is
itself a nontrivial capability (\citeproc{ref-huang2023metatool}{Huang
et al. 2023}).

This paper studies a narrower question:

\begin{quote}
Given a user query, an obligation profile, and a declared registry of
runtime modules, can we compile the minimal harness that satisfies the
external obligations required for a warranted answer or action?
\end{quote}

GapHarness separates the problem into two stages. First, a profiler
infers an obligation profile: obligations, required capabilities, output
contracts, risk level, and unsupported or clarification conditions.
Second, an exact compiler searches a declared module registry for the
lowest-cost module subset that covers the obligations and capabilities.
This separation makes failure interpretable. If a registry does not
contain an affordance, the system should return unsupported rather than
hallucinating support. If the profiler over- or under-infers
obligations, the verifier should expose measurable over-harnessing,
under-harnessing, or wrong-harnessing.

The goal is not to build another general-purpose agent framework. The
goal is to make harness synthesis measurable and auditable. Minimality
is relative to a declared obligation ontology, module registry,
dependency model, and cost function. This is a controlled technical
claim, not a claim of universal agent optimality.

\subsection{Contributions}\label{contributions}

\begin{enumerate}
\def\labelenumi{\arabic{enumi}.}
\tightlist
\item
  We formulate agent harnessing as obligation inference plus minimal
  runtime compilation.
\item
  We implement GapHarness, an API-only compiler from obligations to
  minimal harness modules.
\item
  We introduce minimal-sufficiency metrics: over-harnessing,
  under-harnessing, wrong-harnessing, minimality regret, and redundancy.
\item
  We evaluate GapHarness on GapBench-1000, LLM-inferred profiler
  settings, registry-guarded calibration, transfer subsets, and
  anti-circularity stress tests.
\end{enumerate}

\section{2. Related Work}\label{related-work}

\subsection{Tool-using LLM agents}\label{tool-using-llm-agents}

ReAct interleaves model reasoning traces with environment actions
(\citeproc{ref-yao2023react}{S. Yao et al. 2023}). Toolformer trains
models to decide when and how to call APIs
(\citeproc{ref-schick2023toolformer}{Schick et al. 2023}), while Gorilla
and ToolLLM focus on API selection, API calling, and large tool-use
corpora (\citeproc{ref-patil2023gorilla}{Patil et al. 2023};
\citeproc{ref-qin2024toolllm}{Qin et al. 2024}). MetaTool asks whether
models know when to use tools and which tools to choose
(\citeproc{ref-huang2023metatool}{Huang et al. 2023}). GapHarness
differs by placing obligation identification before tool or module
selection. The compiler does not decide that a query needs ``a search
tool''; it decides whether the profile imposes Observation,
Verification, Execution, State, Action, or Control obligations, then
compiles registry modules that cover those obligations.

\subsection{Retrieval, execution, and
verification}\label{retrieval-execution-and-verification}

Observation and Verification obligations are motivated by
retrieval-augmented generation and provenance-aware answering
(\citeproc{ref-lewis2020rag}{Lewis et al. 2020}). Execution obligations
are motivated by program-aided reasoning, where a language model
delegates deterministic computation to a runtime
(\citeproc{ref-gao2022pal}{Gao et al. 2022}). GapHarness treats these as
module affordances that can be required, absent, or redundant. Its
verifier checks whether the chosen harness covers declared obligations
and capabilities; it does not claim semantic answer correctness for
arbitrary open-world questions.

\subsection{Agent benchmarks and transfer
settings}\label{agent-benchmarks-and-transfer-settings}

AgentBench evaluates LLMs as agents across interactive environments
(\citeproc{ref-liu2024agentbench}{Liu et al. 2024}). GAIA evaluates
general assistant questions requiring reasoning, multimodality, web
browsing, and tool use (\citeproc{ref-mialon2024gaia}{Mialon et al.
2024}). Terminal-Bench evaluates hard terminal-environment tasks
(\citeproc{ref-merrill2026terminalbench}{Merrill et al. 2026}).
WildToolBench and MCP-Bench emphasize realistic tool-use behavior and
cross-tool coordination (\citeproc{ref-yu2026wildtoolbench}{Yu et al.
2026}; \citeproc{ref-wang2025mcpbench}{Wang et al. 2025}). GapHarness
uses GAIA-derived and terminal-derived transfer artifacts only as
obligation-transfer or scaffold evidence. It does not claim full GAIA
answer-level solving or Terminal-Bench solving.

\subsection{Workflow and harness
systems}\label{workflow-and-harness-systems}

AutoFlow, AFlow, and WorFBench study workflow generation or workflow
optimization for agentic systems (\citeproc{ref-li2024autoflow}{Li et
al. 2024}; \citeproc{ref-zhang2025aflow}{Zhang et al. 2025};
\citeproc{ref-qiao2025worfbench}{Qiao et al. 2025}). Recent harness work
treats the execution layer itself as an object of study, including
automated code harness synthesis, natural-language harnesses, and
benchmarked harness effects (\citeproc{ref-lou2026autoharness}{Lou et
al. 2026}; \citeproc{ref-pan2026naturalharness}{Pan et al. 2026};
\citeproc{ref-yao2026harnessbench}{Y. Yao et al. 2026}). GapHarness is
complementary: it studies the compiler-like layer that maps obligations
to the smallest declared runtime harness rather than optimizing a
complete workflow or measuring complete model-harness configurations.

\section{3. Problem Formulation}\label{problem-formulation}

Let a user query be represented by an obligation profile:

\begin{itemize}
\tightlist
\item
  obligations: a subset of Observation, Execution, State, Action,
  Control, Verification;
\item
  required capabilities: lower-level affordances such as
  \texttt{evidence\_sources}, \texttt{execution},
  \texttt{workspace\_inspection}, \texttt{durable\_state},
  \texttt{sandbox\_action}, \texttt{permission},
  \texttt{contract\_check}, or \texttt{source\_spans};
\item
  output contract and risk metadata;
\item
  clarification or unsupported conditions.
\end{itemize}

Let a registry contain modules. Each module declares:

\begin{itemize}
\tightlist
\item
  obligations it provides;
\item
  capabilities it provides;
\item
  dependency requirements over obligations and capabilities;
\item
  cost;
\item
  verifier metadata.
\end{itemize}

The compiler searches for the lowest-cost subset of modules whose
provided obligations and capabilities cover the profile and whose
dependencies are internally satisfied. If no subset covers the profile,
the compiler returns unsupported. If the profile requires clarification,
the compiler returns clarify.

This setup yields three useful failure modes:

\begin{itemize}
\tightlist
\item
  under-harnessing: a supported task lacks required obligation or
  capability coverage;
\item
  over-harnessing: the selected harness costs more than the oracle
  minimal harness;
\item
  wrong-harnessing: a nonempty harness is selected but verifier failures
  remain.
\end{itemize}

Minimality is not absolute. It is relative to the declared registry and
cost function.

\section{4. GapHarness}\label{gapharness}

GapHarness contains five components.

The profiler maps a query to an obligation profile. The current system
supports gold profiles, heuristic profiles, LLM-inferred profiles, and a
registry-guarded LLM profile variant. The gold profiler is used to
isolate compiler behavior. The LLM profiler is used to measure practical
inference quality.

The registry declares available modules. The MVP registry includes web
retrieval, source-span checking, Python execution, execution-log
checking, workspace inspection, state storage, sandbox file editing,
permission gating, contract verification, and trace recording.

The compiler performs exact search over registry subsets and selects the
lowest-cost valid harness. It returns direct-answer, supported, clarify,
or unsupported harness status.

The executor is a deterministic sandbox/mock runtime. It records traces
but does not perform irreversible external side effects.

The verifier checks sufficiency against benchmark gold labels and
computes drop-one minimality reports.

\begin{figure}
\centering
\includegraphics[width=0.95\linewidth,height=\textheight,keepaspectratio]{paper/figures/figure1_pipeline_print.png}
\caption{GapHarness pipeline. The figure shows the obligation profiler,
declared registry, exact compiler, sandbox executor, trace recorder, and
verifier path. It proves system decomposition; it does not prove
open-world task solving.}\label{fig:pipeline}
\end{figure}

\section{5. Benchmarks and Metrics}\label{benchmarks-and-metrics}

\subsection{GapBench}\label{gapbench}

GapBench v1.0 is a 1000-row controlled benchmark designed to isolate
obligation coverage and minimal harness compilation. It contains
human-audited labels for obligations, required capabilities, oracle
minimal harnesses, expected status, success checker, risk level, and
provenance. The benchmark includes dev200 and test800 splits.

GapBench is controlled and factorial. It is not a complete real-world
benchmark. Its purpose is to make the compiler and verifier problem
measurable.

\subsection{Transfer and review sets}\label{transfer-and-review-sets}

GAIA-Transfer v1.0 contains 200 GAIA-derived obligation-transfer rows.
It evaluates obligation assignment and harness selection only; it is not
a full GAIA answer-level benchmark.

GapBench-Natural v1.0 contains 200 naturalized prompts inherited from
GapBench rows. It is a for-review artifact until naturalized prompts are
audited.

Terminal-Bench-obligation50 is an execution-heavy scaffold derived from
public Terminal-Bench task instructions. It is appendix material with
labels pending audit, not a Terminal-Bench solving result.

\subsection{Metrics}\label{metrics}

We report success, average cost, oracle cost, minimality regret,
over-harnessing, under-harnessing, wrong-harnessing, and redundancy.
Success is deterministic verifier pass against expected status and gold
obligation/capability coverage.

\section{6. Experiments}\label{experiments}

\subsection{Gold-obligation compiler
evaluation}\label{gold-obligation-compiler-evaluation}

Table 1 reports GapBench-1000 under gold obligations. GapHarness reaches
1.00 success, average cost 3.67, and 0.00 regret, matching oracle
minimal. Direct answering reaches 0.20 success with 0.74
under-harnessing. Always-full reaches 0.94 success but over-harnesses at
0.94 with average cost 16.00. Tool Router and Difficulty Router show
mixed under- and wrong-harnessing failures.

Table 1 proves that the exact compiler-verifier path matches oracle
minimal harnesses under human-audited gold obligations. It does not
prove open-world LLM profiling or answer-level solving.

\begin{longtable}[]{@{}
  >{\raggedright\arraybackslash}p{(\linewidth - 18\tabcolsep) * \real{0.0769}}
  >{\raggedleft\arraybackslash}p{(\linewidth - 18\tabcolsep) * \real{0.1026}}
  >{\raggedleft\arraybackslash}p{(\linewidth - 18\tabcolsep) * \real{0.1026}}
  >{\raggedleft\arraybackslash}p{(\linewidth - 18\tabcolsep) * \real{0.1026}}
  >{\raggedleft\arraybackslash}p{(\linewidth - 18\tabcolsep) * \real{0.1026}}
  >{\raggedleft\arraybackslash}p{(\linewidth - 18\tabcolsep) * \real{0.1026}}
  >{\raggedleft\arraybackslash}p{(\linewidth - 18\tabcolsep) * \real{0.1026}}
  >{\raggedleft\arraybackslash}p{(\linewidth - 18\tabcolsep) * \real{0.1026}}
  >{\raggedleft\arraybackslash}p{(\linewidth - 18\tabcolsep) * \real{0.1026}}
  >{\raggedleft\arraybackslash}p{(\linewidth - 18\tabcolsep) * \real{0.1026}}@{}}
\toprule\noalign{}
\begin{minipage}[b]{\linewidth}\raggedright
System
\end{minipage} & \begin{minipage}[b]{\linewidth}\raggedleft
N
\end{minipage} & \begin{minipage}[b]{\linewidth}\raggedleft
Success
\end{minipage} & \begin{minipage}[b]{\linewidth}\raggedleft
Avg Cost
\end{minipage} & \begin{minipage}[b]{\linewidth}\raggedleft
Oracle Cost
\end{minipage} & \begin{minipage}[b]{\linewidth}\raggedleft
Regret
\end{minipage} & \begin{minipage}[b]{\linewidth}\raggedleft
Over
\end{minipage} & \begin{minipage}[b]{\linewidth}\raggedleft
Under
\end{minipage} & \begin{minipage}[b]{\linewidth}\raggedleft
Wrong
\end{minipage} & \begin{minipage}[b]{\linewidth}\raggedleft
Redundancy
\end{minipage} \\
\midrule\noalign{}
\endhead
\bottomrule\noalign{}
\endlastfoot
direct & 1000 & 0.20 & 0.00 & 3.67 & -3.67 & 0.00 & 0.74 & 0.00 &
0.00 \\
tool\_router & 1000 & 0.34 & 2.10 & 3.67 & -1.57 & 0.11 & 0.60 & 0.42 &
0.06 \\
difficulty\_router & 1000 & 0.43 & 3.46 & 3.67 & -0.21 & 0.28 & 0.51 &
0.16 & 0.14 \\
always\_full & 1000 & 0.94 & 16.00 & 3.67 & 12.33 & 0.94 & 0.00 & 0.00 &
0.51 \\
gapharness & 1000 & 1.00 & 3.67 & 3.67 & 0.00 & 0.00 & 0.00 & 0.00 &
0.00 \\
oracle\_minimal & 1000 & 1.00 & 3.67 & 3.67 & 0.00 & 0.00 & 0.00 & 0.00
& 0.00 \\
\end{longtable}

\begin{figure}
\centering
\includegraphics[width=0.95\linewidth,height=\textheight,keepaspectratio]{paper/figures/figure2_cost_success_frontier.png}
\caption{Cost-success frontier. The figure shows that Always-full buys
high success with high cost, while gold GapHarness reaches oracle
minimal cost. It does not prove that LLM profiling is fully
calibrated.}\label{fig:frontier}
\end{figure}

\begin{figure}
\centering
\includegraphics[width=0.95\linewidth,height=\textheight,keepaspectratio]{paper/figures/figure3_over_under_wrong_bars.png}
\caption{Over, under, and wrong harnessing rates. The figure separates
baseline failure modes. It does not claim these rates transfer unchanged
to all naturalistic tasks.}\label{fig:ouw}
\end{figure}

\subsection{LLM-inferred obligation
profiling}\label{llm-inferred-obligation-profiling}

Phase 2B evaluates whether an LLM profiler can infer obligations well
enough for practical harness synthesis. On held-out test800, the
selected LLM GapHarness reaches 0.89 success at 3.59 average cost,
outperforming direct, tool-router, and difficulty-router baselines. It
remains below the gold oracle, with under-harnessing at 0.09.

Table 2 proves that inferred obligations are useful against simple
baselines. It does not prove that the profiler is fully calibrated.

\begin{longtable}[]{@{}
  >{\raggedright\arraybackslash}p{(\linewidth - 18\tabcolsep) * \real{0.0769}}
  >{\raggedleft\arraybackslash}p{(\linewidth - 18\tabcolsep) * \real{0.1026}}
  >{\raggedleft\arraybackslash}p{(\linewidth - 18\tabcolsep) * \real{0.1026}}
  >{\raggedleft\arraybackslash}p{(\linewidth - 18\tabcolsep) * \real{0.1026}}
  >{\raggedleft\arraybackslash}p{(\linewidth - 18\tabcolsep) * \real{0.1026}}
  >{\raggedleft\arraybackslash}p{(\linewidth - 18\tabcolsep) * \real{0.1026}}
  >{\raggedleft\arraybackslash}p{(\linewidth - 18\tabcolsep) * \real{0.1026}}
  >{\raggedleft\arraybackslash}p{(\linewidth - 18\tabcolsep) * \real{0.1026}}
  >{\raggedleft\arraybackslash}p{(\linewidth - 18\tabcolsep) * \real{0.1026}}
  >{\raggedleft\arraybackslash}p{(\linewidth - 18\tabcolsep) * \real{0.1026}}@{}}
\toprule\noalign{}
\begin{minipage}[b]{\linewidth}\raggedright
System
\end{minipage} & \begin{minipage}[b]{\linewidth}\raggedleft
N
\end{minipage} & \begin{minipage}[b]{\linewidth}\raggedleft
Success
\end{minipage} & \begin{minipage}[b]{\linewidth}\raggedleft
Avg Cost
\end{minipage} & \begin{minipage}[b]{\linewidth}\raggedleft
Oracle Cost
\end{minipage} & \begin{minipage}[b]{\linewidth}\raggedleft
Regret
\end{minipage} & \begin{minipage}[b]{\linewidth}\raggedleft
Over
\end{minipage} & \begin{minipage}[b]{\linewidth}\raggedleft
Under
\end{minipage} & \begin{minipage}[b]{\linewidth}\raggedleft
Wrong
\end{minipage} & \begin{minipage}[b]{\linewidth}\raggedleft
Redundancy
\end{minipage} \\
\midrule\noalign{}
\endhead
\bottomrule\noalign{}
\endlastfoot
direct & 800 & 0.20 & 0.00 & 3.69 & -3.69 & 0.00 & 0.74 & 0.00 & 0.00 \\
tool\_router & 800 & 0.32 & 1.96 & 3.69 & -1.72 & 0.09 & 0.62 & 0.43 &
0.06 \\
difficulty\_router & 800 & 0.41 & 3.22 & 3.69 & -0.47 & 0.26 & 0.53 &
0.15 & 0.13 \\
always\_full & 800 & 0.94 & 16.00 & 3.69 & 12.31 & 0.94 & 0.00 & 0.00 &
0.51 \\
gold\_oracle\_gap\_harness & 800 & 1.00 & 3.69 & 3.69 & 0.00 & 0.00 &
0.00 & 0.00 & 0.00 \\
selected\_llm\_gap\_harness & 800 & 0.89 & 3.59 & 3.69 & -0.09 & 0.14 &
0.09 & 0.01 & 0.05 \\
\end{longtable}

\subsection{Registry-guarded
calibration}\label{registry-guarded-calibration}

Phase 2B exposed a systematic failure mode: sandbox/mock/local actions
were sometimes lowered into unsupported
\texttt{real\_world\_side\_effect}. Phase 2C adds a deterministic
registry guard that preserves real external side-effect boundaries but
removes \texttt{real\_world\_side\_effect} when the query explicitly
limits action to sandbox/mock/local scope.

On held-out test800, registry guarding improves success from 0.89 to
0.94 and reduces under-harnessing from 0.09 to 0.03. Unsupported false
positives fall from 56 to 12. The tradeoff is higher average cost and
regret: cost rises from 3.59 to 3.98 and regret from -0.09 to 0.30. This
supports a narrow calibration claim: declared registry boundaries can
repair a specific profiler failure mode.

Table 3 and Figure 4 prove that registry guarding repairs this
sandbox-action unsupported false-positive pattern. They do not prove a
general solution to all profiling failures.

\begin{longtable}[]{@{}
  >{\raggedright\arraybackslash}p{(\linewidth - 18\tabcolsep) * \real{0.0789}}
  >{\raggedright\arraybackslash}p{(\linewidth - 18\tabcolsep) * \real{0.0789}}
  >{\raggedleft\arraybackslash}p{(\linewidth - 18\tabcolsep) * \real{0.1053}}
  >{\raggedleft\arraybackslash}p{(\linewidth - 18\tabcolsep) * \real{0.1053}}
  >{\raggedleft\arraybackslash}p{(\linewidth - 18\tabcolsep) * \real{0.1053}}
  >{\raggedleft\arraybackslash}p{(\linewidth - 18\tabcolsep) * \real{0.1053}}
  >{\raggedleft\arraybackslash}p{(\linewidth - 18\tabcolsep) * \real{0.1053}}
  >{\raggedleft\arraybackslash}p{(\linewidth - 18\tabcolsep) * \real{0.1053}}
  >{\raggedleft\arraybackslash}p{(\linewidth - 18\tabcolsep) * \real{0.1053}}
  >{\raggedleft\arraybackslash}p{(\linewidth - 18\tabcolsep) * \real{0.1053}}@{}}
\toprule\noalign{}
\begin{minipage}[b]{\linewidth}\raggedright
Split
\end{minipage} & \begin{minipage}[b]{\linewidth}\raggedright
Profiler
\end{minipage} & \begin{minipage}[b]{\linewidth}\raggedleft
N
\end{minipage} & \begin{minipage}[b]{\linewidth}\raggedleft
Success
\end{minipage} & \begin{minipage}[b]{\linewidth}\raggedleft
Avg Cost
\end{minipage} & \begin{minipage}[b]{\linewidth}\raggedleft
Regret
\end{minipage} & \begin{minipage}[b]{\linewidth}\raggedleft
Over
\end{minipage} & \begin{minipage}[b]{\linewidth}\raggedleft
Under
\end{minipage} & \begin{minipage}[b]{\linewidth}\raggedleft
Wrong
\end{minipage} & \begin{minipage}[b]{\linewidth}\raggedleft
Unsupported FP
\end{minipage} \\
\midrule\noalign{}
\endhead
\bottomrule\noalign{}
\endlastfoot
dev200 & Phase 2B llm\_single & 200 & 0.92 & 3.68 & 0.06 & 0.19 & 0.08 &
0.00 & 14 \\
dev200 & Phase 2C registry\_guarded & 200 & 0.97 & 4.02 & 0.40 & 0.20 &
0.03 & 0.00 & 4 \\
test800 & Phase 2B selected llm\_single & 800 & 0.89 & 3.59 & -0.09 &
0.14 & 0.09 & 0.01 & 56 \\
test800 & Phase 2C registry\_guarded & 800 & 0.94 & 3.98 & 0.30 & 0.15 &
0.03 & 0.01 & 12 \\
\end{longtable}

\begin{figure}
\centering
\includegraphics[width=0.9\linewidth,height=\textheight,keepaspectratio]{paper/figures/figure4_registry_guard_unsupported_fp_reduction.png}
\caption{Registry guard correction. The figure shows reduced unsupported
false positives and under-harnessing after a registry guard. It does not
prove that registry guarding solves GAIA or all transfer
settings.}\label{fig:guard}
\end{figure}

\subsection{Stress tests and negative
controls}\label{stress-tests-and-negative-controls}

Registry perturbation removes six key modules one at a time. For each
relevant 60-task subset, base registry success is 1.00 and perturbed
registry success is 0.00. The dominant missing capabilities match the
removed module: \texttt{execution}, \texttt{source\_spans},
\texttt{permission}, \texttt{diff}, \texttt{evidence\_sources}, and
\texttt{contract\_check}. This shows that GapHarness does not silently
claim support when affordances are absent.

Gold label permutation corrupts 200 supported profiles, while the
verifier still checks original human-audited labels. Correct labels
yield 1.00 success. Permuted labels reduce success to 0.17, raise
under-harnessing to 0.83, and raise wrong-harnessing to 0.79. The
permutation generator changes obligations or required capabilities for
all 200 corrupted profiles. This is an anti-circularity test, not a
realistic noise model.

Negative controls evaluate pure-language and tool-bait prompts.
GapHarness gold, LLM, and registry-guarded variants all achieve 1.00
success, 0.00 average cost, and 0.00 over-harnessing. Always-full
over-harnesses both categories; Tool Router and Difficulty Router
over-harness tool-bait prompts at 0.51.

Table 4 proves three anti-circularity properties: registry affordances
matter, labels matter, and tool-like words alone do not trigger
GapHarness over-harnessing. It does not model all realistic label noise
or all adversarial prompts.

\begin{longtable}[]{@{}
  >{\raggedright\arraybackslash}p{(\linewidth - 10\tabcolsep) * \real{0.1364}}
  >{\raggedleft\arraybackslash}p{(\linewidth - 10\tabcolsep) * \real{0.1818}}
  >{\raggedleft\arraybackslash}p{(\linewidth - 10\tabcolsep) * \real{0.1818}}
  >{\raggedleft\arraybackslash}p{(\linewidth - 10\tabcolsep) * \real{0.1818}}
  >{\raggedleft\arraybackslash}p{(\linewidth - 10\tabcolsep) * \real{0.1818}}
  >{\raggedright\arraybackslash}p{(\linewidth - 10\tabcolsep) * \real{0.1364}}@{}}
\toprule\noalign{}
\begin{minipage}[b]{\linewidth}\raggedright
Removed module
\end{minipage} & \begin{minipage}[b]{\linewidth}\raggedleft
Base
\end{minipage} & \begin{minipage}[b]{\linewidth}\raggedleft
Perturbed
\end{minipage} & \begin{minipage}[b]{\linewidth}\raggedleft
Unsupported
\end{minipage} & \begin{minipage}[b]{\linewidth}\raggedleft
Under
\end{minipage} & \begin{minipage}[b]{\linewidth}\raggedright
Missing capability
\end{minipage} \\
\midrule\noalign{}
\endhead
\bottomrule\noalign{}
\endlastfoot
python\_executor & 1.00 & 0.00 & 1.00 & 1.00 & execution \\
source\_span\_checker & 1.00 & 0.00 & 1.00 & 1.00 & source\_spans \\
permission\_gate & 1.00 & 0.00 & 1.00 & 1.00 & permission \\
sandbox\_file\_editor & 1.00 & 0.00 & 1.00 & 1.00 & diff \\
web\_retrieval & 1.00 & 0.00 & 1.00 & 1.00 & evidence\_sources \\
contract\_verifier & 1.00 & 0.00 & 1.00 & 1.00 & contract\_check \\
\end{longtable}

\begin{longtable}[]{@{}
  >{\raggedright\arraybackslash}p{(\linewidth - 14\tabcolsep) * \real{0.0968}}
  >{\raggedleft\arraybackslash}p{(\linewidth - 14\tabcolsep) * \real{0.1290}}
  >{\raggedleft\arraybackslash}p{(\linewidth - 14\tabcolsep) * \real{0.1290}}
  >{\raggedleft\arraybackslash}p{(\linewidth - 14\tabcolsep) * \real{0.1290}}
  >{\raggedleft\arraybackslash}p{(\linewidth - 14\tabcolsep) * \real{0.1290}}
  >{\raggedleft\arraybackslash}p{(\linewidth - 14\tabcolsep) * \real{0.1290}}
  >{\raggedleft\arraybackslash}p{(\linewidth - 14\tabcolsep) * \real{0.1290}}
  >{\raggedleft\arraybackslash}p{(\linewidth - 14\tabcolsep) * \real{0.1290}}@{}}
\toprule\noalign{}
\begin{minipage}[b]{\linewidth}\raggedright
Condition
\end{minipage} & \begin{minipage}[b]{\linewidth}\raggedleft
N
\end{minipage} & \begin{minipage}[b]{\linewidth}\raggedleft
Success
\end{minipage} & \begin{minipage}[b]{\linewidth}\raggedleft
Regret
\end{minipage} & \begin{minipage}[b]{\linewidth}\raggedleft
Over
\end{minipage} & \begin{minipage}[b]{\linewidth}\raggedleft
Under
\end{minipage} & \begin{minipage}[b]{\linewidth}\raggedleft
Wrong
\end{minipage} & \begin{minipage}[b]{\linewidth}\raggedleft
Verifier Fail
\end{minipage} \\
\midrule\noalign{}
\endhead
\bottomrule\noalign{}
\endlastfoot
correct gold & 200 & 1.00 & 0.00 & 0.00 & 0.00 & 0.00 & 0.00 \\
permuted labels & 200 & 0.17 & 0.24 & 0.55 & 0.83 & 0.79 & 0.83 \\
\end{longtable}

\begin{longtable}[]{@{}llrrrr@{}}
\toprule\noalign{}
Category & System & N & Success & Cost & Over \\
\midrule\noalign{}
\endhead
\bottomrule\noalign{}
\endlastfoot
pure-lang & GH gold & 100 & 1.00 & 0.00 & 0.00 \\
pure-lang & GH LLM & 100 & 1.00 & 0.00 & 0.00 \\
pure-lang & GH guarded & 100 & 1.00 & 0.00 & 0.00 \\
pure-lang & Always-full & 100 & 1.00 & 16.00 & 1.00 \\
tool-bait & GH gold & 100 & 1.00 & 0.00 & 0.00 \\
tool-bait & GH LLM & 100 & 1.00 & 0.00 & 0.00 \\
tool-bait & GH guarded & 100 & 1.00 & 0.00 & 0.00 \\
tool-bait & Tool Router & 100 & 1.00 & 1.26 & 0.51 \\
tool-bait & Difficulty & 100 & 1.00 & 1.22 & 0.51 \\
tool-bait & Always-full & 100 & 1.00 & 16.00 & 1.00 \\
\end{longtable}

\subsection{Transfer and boundary
results}\label{transfer-and-boundary-results}

GAIA-Transfer gold smoke reaches 1.00 success under gold obligation
labels. This confirms that the compiler can process the transfer labels,
not that the system solves GAIA.

GAIA registry-guarded reaches 0.56 success with high over-harnessing and
under-harnessing. This negative result is useful: the GapBench registry
guard targets sandbox/mock action calibration, while GAIA failures are
dominated by multimodal, file, evidence, and state-boundary mismatches.

GapBench-Natural review smoke reaches 1.00 success under inherited
labels, but should not be treated as final until naturalized prompts are
audited.

Table 5 proves that the representation can be applied to transfer
artifacts and exposes boundary failures. It does not claim full GAIA
solving, Terminal-Bench solving, or final natural-prompt validation.

\begin{longtable}[]{@{}
  >{\raggedright\arraybackslash}p{(\linewidth - 16\tabcolsep) * \real{0.0882}}
  >{\raggedright\arraybackslash}p{(\linewidth - 16\tabcolsep) * \real{0.0882}}
  >{\raggedleft\arraybackslash}p{(\linewidth - 16\tabcolsep) * \real{0.1176}}
  >{\raggedleft\arraybackslash}p{(\linewidth - 16\tabcolsep) * \real{0.1176}}
  >{\raggedleft\arraybackslash}p{(\linewidth - 16\tabcolsep) * \real{0.1176}}
  >{\raggedleft\arraybackslash}p{(\linewidth - 16\tabcolsep) * \real{0.1176}}
  >{\raggedleft\arraybackslash}p{(\linewidth - 16\tabcolsep) * \real{0.1176}}
  >{\raggedleft\arraybackslash}p{(\linewidth - 16\tabcolsep) * \real{0.1176}}
  >{\raggedleft\arraybackslash}p{(\linewidth - 16\tabcolsep) * \real{0.1176}}@{}}
\toprule\noalign{}
\begin{minipage}[b]{\linewidth}\raggedright
Result
\end{minipage} & \begin{minipage}[b]{\linewidth}\raggedright
Identity
\end{minipage} & \begin{minipage}[b]{\linewidth}\raggedleft
N
\end{minipage} & \begin{minipage}[b]{\linewidth}\raggedleft
Success
\end{minipage} & \begin{minipage}[b]{\linewidth}\raggedleft
Cost
\end{minipage} & \begin{minipage}[b]{\linewidth}\raggedleft
Oracle
\end{minipage} & \begin{minipage}[b]{\linewidth}\raggedleft
Regret
\end{minipage} & \begin{minipage}[b]{\linewidth}\raggedleft
Over
\end{minipage} & \begin{minipage}[b]{\linewidth}\raggedleft
Under
\end{minipage} \\
\midrule\noalign{}
\endhead
\bottomrule\noalign{}
\endlastfoot
GAIA-Transfer gold & transfer smoke & 200 & 1.00 & 1.48 & 1.48 & 0.00 &
0.00 & 0.00 \\
GAIA guarded & limitation & 200 & 0.56 & 5.56 & 1.48 & 4.08 & 0.89 &
0.44 \\
GapBench-Natural & review smoke & 200 & 1.00 & 2.83 & 2.83 & 0.00 & 0.00
& 0.00 \\
Terminal obligation50 & appendix scaffold & 50 & - & - & - & - & - &
- \\
\end{longtable}

Boundary notes: GAIA-Transfer evaluates obligation labels only, not GAIA
answer solving. GAIA guarded is a limitation result dominated by
multimodal, file, evidence, and state-boundary gaps. GapBench-Natural
inherits labels pending natural-prompt audit. Terminal obligation50 is
not a Terminal-Bench solving result.

\section{7. Analysis}\label{analysis}

The experimental story separates three layers.

First, under gold obligations, minimal harness compilation works. The
compiler matches oracle minimal cost and produces interpretable failures
when the registry lacks affordances.

Second, LLM profiling is the main practical bottleneck. LLM-inferred
obligations outperform simple routers but can under-harness or
over-harness depending on calibration. Registry guarding shows that some
errors are not random; they can be tied to declared affordance
boundaries.

Third, stress tests reduce circularity concerns. Registry perturbation
shows the system depends on declared affordances. Label permutation
shows labels carry semantic force. Negative controls show the system is
not merely responding to tool-like words.

These results support a controlled claim: obligation-first minimal
harness synthesis is a useful way to separate profiler errors from
compiler/runtime errors.

\section{8. Limitations}\label{limitations}

GapHarness does not solve full GAIA. GAIA-Transfer is
obligation-transfer only and does not evaluate answer-level accuracy
against GAIA final answers.

GapHarness does not solve Terminal-Bench. Terminal-Bench-obligation50 is
a scaffold derived from task instructions, with labels pending audit.

The executor is a deterministic sandbox/mock runtime. It does not
perform irreversible file edits, real API calls, payments, emails,
deployments, or production changes.

GapBench is controlled and factorial. This is a strength for isolating
compiler behavior, but it is not a substitute for large-scale
naturalistic agent benchmarks.

The LLM profiler remains imperfect. Registry guarding improves one
systematic failure mode, but GAIA-Transfer shows that broader transfer
remains limited.

Minimality is registry-relative. Different module registries, cost
functions, or dependency models would change the compiled minimal
harness.

\section{9. Conclusion}\label{conclusion}

GapHarness reframes agent harnessing as two separable problems:
identifying the external obligations imposed by a query, and compiling
the smallest runtime system that satisfies those obligations. Under
human-audited gold obligations, the compiler matches oracle minimal
harnesses on GapBench-1000. With LLM-inferred obligations, GapHarness
improves over direct and router baselines but exposes calibration
tradeoffs. Registry guarding repairs a systematic unsupported
false-positive failure. Stress tests show that success depends on
declared registry affordances and meaningful obligation labels, while
negative controls show that GapHarness avoids tool-bait over-harnessing.

The current evidence supports a workshop/arXiv technical report claim:
obligation-first minimal harness synthesis provides a measurable,
auditable layer for API-only LLM agents. The next step is not to claim
broad open-world solving, but to expand audited transfer benchmarks and
improve LLM obligation profiling under the same evidence boundaries.

\section*{References}\label{references}
\addcontentsline{toc}{section}{References}

\protect\phantomsection\label{refs}
\begin{CSLReferences}{1}{0}
\bibitem[\citeproctext]{ref-gao2022pal}
Gao, Luyu, Aman Madaan, Shuyan Zhou, Uri Alon, Pengfei Liu, Yiming Yang,
Jamie Callan, and Graham Neubig. 2022. {``PAL: Program-Aided Language
Models.''} \emph{arXiv Preprint arXiv:2211.10435}.
\url{https://arxiv.org/abs/2211.10435}.

\bibitem[\citeproctext]{ref-huang2023metatool}
Huang, Yue, Jiawen Shi, Yuan Li, Chenrui Fan, Siyuan Wu, Qihui Zhang,
Yixin Liu, et al. 2023. {``MetaTool Benchmark for Large Language Models:
Deciding Whether to Use Tools and Which to Use.''} \emph{arXiv Preprint
arXiv:2310.03128}. \url{https://arxiv.org/abs/2310.03128}.

\bibitem[\citeproctext]{ref-lewis2020rag}
Lewis, Patrick, Ethan Perez, Aleksandra Piktus, Fabio Petroni, Vladimir
Karpukhin, Naman Goyal, Heinrich Küttler, et al. 2020.
{``Retrieval-Augmented Generation for Knowledge-Intensive NLP Tasks.''}
In \emph{Advances in Neural Information Processing Systems}.
\url{https://arxiv.org/abs/2005.11401}.

\bibitem[\citeproctext]{ref-li2024autoflow}
Li, Zelong, Shuyuan Xu, Kai Mei, Wenyue Hua, Balaji Rama, Om Raheja, Hao
Wang, He Zhu, and Yongfeng Zhang. 2024. {``AutoFlow: Automated Workflow
Generation for Large Language Model Agents.''} \emph{arXiv Preprint
arXiv:2407.12821}. \url{https://arxiv.org/abs/2407.12821}.

\bibitem[\citeproctext]{ref-liu2024agentbench}
Liu, Xiao, Hao Yu, Hanchen Zhang, Yifan Xu, Xuanyu Lei, Hanyu Lai, Yu
Gu, et al. 2024. {``AgentBench: Evaluating LLMs as Agents.''} In
\emph{International Conference on Learning Representations}.
\url{https://openreview.net/forum?id=zAdUB0aCTQ}.

\bibitem[\citeproctext]{ref-lou2026autoharness}
Lou, Xinghua, Miguel Lázaro-Gredilla, Antoine Dedieu, Carter Wendelken,
Wolfgang Lehrach, and Kevin P. Murphy. 2026. {``AutoHarness: Improving
LLM Agents by Automatically Synthesizing a Code Harness.''} \emph{arXiv
Preprint arXiv:2603.03329}. \url{https://arxiv.org/abs/2603.03329}.

\bibitem[\citeproctext]{ref-merrill2026terminalbench}
Merrill, Mike A., Alexander G. Shaw, Nicholas Carlini, Boxuan Li, Harsh
Raj, Ivan Bercovich, Lin Shi, et al. 2026. {``Terminal-Bench:
Benchmarking Agents on Hard, Realistic Tasks in Command Line
Interfaces.''} \emph{arXiv Preprint arXiv:2601.11868}.
\url{https://arxiv.org/abs/2601.11868}.

\bibitem[\citeproctext]{ref-mialon2024gaia}
Mialon, Grégoire, Clémentine Fourrier, Craig Swift, Thomas Wolf, Yann
LeCun, and Thomas Scialom. 2024. {``GAIA: A Benchmark for General AI
Assistants.''} In \emph{International Conference on Learning
Representations}. \url{https://arxiv.org/abs/2311.12983}.

\bibitem[\citeproctext]{ref-pan2026naturalharness}
Pan, Linyue, Lexiao Zou, Shuo Guo, Jingchen Ni, and Hai-Tao Zheng. 2026.
{``Natural-Language Agent Harnesses.''} \emph{arXiv Preprint
arXiv:2603.25723}. \url{https://arxiv.org/abs/2603.25723}.

\bibitem[\citeproctext]{ref-patil2023gorilla}
Patil, Shishir G., Tianjun Zhang, Xin Wang, and Joseph E. Gonzalez.
2023. {``Gorilla: Large Language Model Connected with Massive APIs.''}
\emph{arXiv Preprint arXiv:2305.15334}.
\url{https://arxiv.org/abs/2305.15334}.

\bibitem[\citeproctext]{ref-qiao2025worfbench}
Qiao, Shuofei, Runnan Fang, Zhisong Qiu, Xiaobin Wang, Ningyu Zhang,
Yong Jiang, Pengjun Xie, Fei Huang, and Huajun Chen. 2025.
{``Benchmarking Agentic Workflow Generation.''} In \emph{International
Conference on Learning Representations}.
\url{https://arxiv.org/abs/2410.07869}.

\bibitem[\citeproctext]{ref-qin2024toolllm}
Qin, Yujia, Shihao Liang, Yining Ye, Kunlun Zhu, Lan Yan, Yaxi Lu,
Yankai Lin, et al. 2024. {``ToolLLM: Facilitating Large Language Models
to Master 16000+ Real-World APIs.''} In \emph{International Conference
on Learning Representations}.
\url{https://openreview.net/forum?id=dHng2O0Jjr}.

\bibitem[\citeproctext]{ref-schick2023toolformer}
Schick, Timo, Jane Dwivedi-Yu, Roberto Dessì, Roberta Raileanu, Maria
Lomeli, Luke Zettlemoyer, Nicola Cancedda, and Thomas Scialom. 2023.
{``Toolformer: Language Models Can Teach Themselves to Use Tools.''} In
\emph{Advances in Neural Information Processing Systems}.
\url{https://arxiv.org/abs/2302.04761}.

\bibitem[\citeproctext]{ref-wang2025mcpbench}
Wang, Zhenting, Qi Chang, Hemani Patel, Shashank Biju, Cheng-En Wu, Quan
Liu, Aolin Ding, et al. 2025. {``MCP-Bench: Benchmarking Tool-Using LLM
Agents with Complex Real-World Tasks via MCP Servers.''} \emph{arXiv
Preprint arXiv:2508.20453}. \url{https://arxiv.org/abs/2508.20453}.

\bibitem[\citeproctext]{ref-yao2023react}
Yao, Shunyu, Jeffrey Zhao, Dian Yu, Nan Du, Izhak Shafran, Karthik
Narasimhan, and Yuan Cao. 2023. {``ReAct: Synergizing Reasoning and
Acting in Language Models.''} In \emph{International Conference on
Learning Representations}. \url{https://arxiv.org/abs/2210.03629}.

\bibitem[\citeproctext]{ref-yao2026harnessbench}
Yao, Yilun, Xinyu Tan, Chao-Hsuan Liu, Yaoming Li, Zhengyang Wang,
Wenhan Yu, Zhewen Tan, et al. 2026. {``Harness-Bench: Measuring Harness
Effects Across Models in Realistic Agent Workflows.''} \emph{arXiv
Preprint arXiv:2605.27922}. \url{https://arxiv.org/abs/2605.27922}.

\bibitem[\citeproctext]{ref-yu2026wildtoolbench}
Yu, Peijie, Wei Liu, Yifan Yang, Jinjian Li, Zelong Zhang, Xiao Feng,
and Feng Zhang. 2026. {``Benchmarking LLM Tool-Use in the Wild.''}
\emph{arXiv Preprint arXiv:2604.06185}.
\url{https://arxiv.org/abs/2604.06185}.

\bibitem[\citeproctext]{ref-zhang2025aflow}
Zhang, Jiayi, Jinyu Xiang, Zhaoyang Yu, Fengwei Teng, Xionghui Chen,
Jiaqi Chen, Mingchen Zhuge, et al. 2025. {``AFlow: Automating Agentic
Workflow Generation.''} In \emph{International Conference on Learning
Representations}. \url{https://openreview.net/forum?id=z5uVAKwmjf}.

\end{CSLReferences}

\end{document}
